% !TEX root = ../main.tex
% =====================================================================
%  VLASTNÍ MAKRA – zde jsou pomocné příkazy, které se používají v main.tex.
%  Běžně je nemusíte měnit.
% =====================================================================

% ---------- šedý text = místo k doplnění ----------
% Použití: \vstup{Text, který mám nahradit}
\newcommand{\vstup}[1]{\textcolor{gray!85!black}{\itshape #1}}

% ---------- obrázek s automatickým číslem a popiskem ----------
% Použití:
%   \obrazek[0.8\textwidth]{obrazek-1.png}{Popisek.}{obr:stitek}
%
% Argumenty:
%   1. (nepovinný) šířka obrázku, výchozí 0.85\textwidth
%   2. název souboru – hledá se v obrazky/ i v kořeni projektu
%   3. popisek pod obrázkem
%   4. štítek pro odkazy, např. \ref{obr:stitek}
%
% Dokud soubor neexistuje, vysází se šedý rámeček se jménem souboru.
\newcommand{\obrazek}[4][0.85\textwidth]{%
  \begin{figure}[H]
    \centering
    \setlength{\fboxsep}{0pt}%
    \IfFileExists{obrazky/#2}{%
      \includegraphics[width=#1,height=0.45\textheight,keepaspectratio]{obrazky/#2}%
    }{%
      \IfFileExists{#2}{%
        \includegraphics[width=#1,height=0.45\textheight,keepaspectratio]{#2}%
      }{%
        \fcolorbox{gray!45}{gray!8}{%
          \begin{minipage}[c][5cm][c]{\dimexpr#1-2\fboxrule\relax}
            \centering\color{gray}\itshape
            místo pro obrázek\\[4pt]
            \texttt{\detokenize{#2}}
          \end{minipage}}%
      }%
    }%
    \caption{#3}%
    \label{#4}%
  \end{figure}%
}
